% !TEX root = ../../main.tex
% ============================================================================
% The E_infty-Topologization Theorem.  Iterated higher-spin Sugawara/Casimir
% tower upgrades the E_3-topological climax to an E_infty-topological ladder.
%
% Canonical upgrade:
%   Virasoro (N=2) -> E_3-top;  W_N -> E_{N+1}-top;  W_infty -> E_infty-top.
%
% Chapter structure:
%   Section 1. Setup: higher-spin stress tower of depth N.
%   Section 2. Theorem (iterated Sugawara/Casimir construction).
%   Section 3. Theorem (E-ladder via topologized Dunn additivity).
%   Section 4. Three specializations (Virasoro, W_N, W_infty).
%   Section 5. Theorem (Operadic spiral and E_infty formal completion).
%   Section 6. Scope clarifications (E_2 to E_3 upgrade, conformal-vector
%     input, non-principal DS, class-M free-PVA scope).
%   Section 7. Climax restatement (3d+infty topological gravity).
% ============================================================================

\section{The \texorpdfstring{$\Einfty$}{E\_infty}-Topologization Theorem}
\label{sec:e-infinity-topologization}

The 3d-gravity climax part (\ref{part:3d-gravity}) culminates, in its
original inscription, with an
$\Ethree$-topological climax: the derived chiral centre of an affine
Kac--Moody algebra, of a principal $\cW$-algebra, or of a freely
generated Poisson vertex algebra, carries an $\Ethree$-topological
algebra structure via the topologization Construction~\ref{constr:topologization}
driven by a single conformal vector $T(z)$. The Sugawara identity
$T = [Q_{\mathrm{CS}}, G]$ purchases one transverse real dimension, the
$\Etwo^{\mathrm{top}}$ upgrade along the holomorphic direction,
and Dunn additivity with the open colour assembles the three-dimensional
topological field theory.

This chapter asserts that $\Ethree$ is the \emph{first rung} of an infinite
ladder. Each higher-spin conformal tensor in the stress tower
$T_{\bullet}(A) = \bigl\{ T^{(2)}, T^{(3)}, T^{(4)}, \ldots \bigr\}$ of a
chiral algebra of the $\cW$-type purchases \emph{one more} transverse real
dimension; a chiral algebra equipped with stress tower of depth
$N$ reaches $\Etopo{N{+}1}$; the $\cW_\infty$ algebra of Bakas--Pope--Romans
and Linshaw\,\cite{Linshaw-Winfty,Bakas-Winfty,BMP-Wsymmetry} reaches the
universal limit $\Einfty$-topological. The $3$d quantum gravity of
Chapter~\ref{ch:3d-gravity} (Virasoro case $N=2$) is then the $\Ethree$
shadow of a $3$d$+\infty$ topological field theory living on the $\cW_\infty$
boundary. The operadic circle of Section~\ref{sec:e_n-circle-holographic}
is, at Platonic level, an infinite bidirectional spiral.

\begin{remark}[Placement in the programme]
\label{rem:e-infinity-placement}
Section~\ref{sec:e-infinity-topologization} is an upgrade of the
$\Ethree$-topological story of
Theorems~\ref{thm:E3-topological-km},
\ref{thm:E3-topological-DS},
\ref{thm:E3-topological-DS-general},
\ref{thm:E3-topological-free-PVA}. Each of those theorems remains valid,
and each becomes the $N=2$ instance of the more general assertion
Theorem~\ref{thm:e-infinity-topologization-ladder}. The chapter
introduces one new hypothesis clause
(Definition~\ref{def:higher-spin-stress-tower}) and one new construction
(Construction~\ref{constr:iterated-topologization}); these collapse to
the existing single-conformal-vector hypothesis and
Construction~\ref{constr:topologization} at depth $N=2$.
\end{remark}

\subsection{Higher-spin stress tower of depth \texorpdfstring{$N$}{N}}
\label{subsec:higher-spin-stress-tower}

The key hypothesis for the $\Einfty$-ladder is the existence, inside the
chiral algebra $A$, of an infinite (or truncated at depth $N$) tower of
conformal tensors of strictly increasing spin, each of which is
\emph{inner Sugawara-like} at non-critical level. The construction of
these tensors in explicit examples is well-documented; see
Bouwknegt--McCarthy--Pilch\,\cite{BMP-Wsymmetry} for a review of the spin
content of $\cW_N$, Linshaw\,\cite{Linshaw-Winfty} for the universal
two-parameter family $\cW_\infty[\mu]$, and
Schoutens--Sevrin--van Nieuwenhuizen\,\cite{Schoutens-higher-spin-Sugawara}
for the higher-spin Casimir construction originating the terminology.

\begin{definition}[Higher-spin stress tower of depth $N$]
\label{def:higher-spin-stress-tower}
\index{stress tower!higher-spin|textbf}
\index{higher-spin stress tower|textbf}
\index{Wbar infinity algebra@$\cW_\infty$-algebra!stress tower}
Let $A$ be a chiral algebra on a smooth curve $X$, at non-critical level
$k \ne -h^\vee$. A \emph{higher-spin stress tower of depth
$N \in \{2, 3, \ldots\} \cup \{\infty\}$ on $A$} is a graded family of
conformal tensors
\begin{equation}\label{eq:stress-tower-family}
T_{\bullet}(A) \;=\;
\bigl\{\, T^{(n)}(z) \in A \ :\ 2 \le n \le N+1 \,\bigr\},
\qquad \deg\, T^{(n)} \;=\; n
\quad\text{(conformal weight)},
\end{equation}
satisfying the following axioms:
\begin{enumerate}[label=(T\arabic*)]
\item\label{tower:quasiprimary}
 \emph{Quasiprimarity}: each $T^{(n)}$ is a quasiprimary field of
 conformal weight $n$ with respect to the Virasoro
 subalgebra generated by $T^{(2)}$.
\item\label{tower:inner-sugawara}
 \emph{Inner Sugawara-like}: each $T^{(n)}$ is a normally ordered
 polynomial of total conformal degree $n$ in a chosen set of
 generating currents of $A$ (Casimir construction); at $n=2$ this
 reduces to the ordinary Sugawara expression
 $T^{(2)} = T_{\mathrm{Sug}} + T_{\mathrm{imp}}$ familiar from the
 proof of Theorem~\ref{thm:E3-topological-DS-general}.
\item\label{tower:truncated-Winfty}
 \emph{Truncated $\cW_{1+\infty}$ brackets}: the OPE algebra of
 $\{T^{(n)}\}_{2 \le n \le N+1}$ is a quotient of the universal
 $\cW_\infty[\mu]$ algebra of Linshaw\,\cite{Linshaw-Winfty} by the
 truncation ideal cutting off spin $>N+1$. In particular the highest
 spin OPE closes on fields of spin $\le N+1$.
\item\label{tower:BRST-primitive}
 \emph{BRST-primitive existence}: there is a $3$d holomorphic-topological
 BV-BRST quantisation $(A^{\mathrm{BV}}_{3d}, Q_{\mathrm{tot}})$ whose
 boundary chiral algebra is $A$ and which admits, for each $2 \le n \le
 N+1$, a spin-$(n{-}1)$ \emph{BRST-primitive antighost current}
 $G^{(n)}(z) \in A^{\mathrm{BV}}_{3d}$ satisfying
 \begin{equation}\label{eq:nth-Sugawara-BRST}
 T^{(n)}(z) \;=\; \bigl[ Q_{\mathrm{tot}},\, G^{(n)}(z) \bigr]
 \qquad\text{on}\ H^{\bullet}(A^{\mathrm{BV}}_{3d}, Q_{\mathrm{tot}}).
 \end{equation}
\item\label{tower:antighost-commutativity}
 \emph{Antighost BRST-commutativity}: the antighosts pairwise
 BRST-commute on cohomology:
 \[
 \bigl[ G^{(n)}(z),\, G^{(m)}(w) \bigr]
 \;=\; Q_{\mathrm{tot}}\text{-exact}
 \qquad \text{for all $2 \le n, m \le N+1$.}
 \]
\end{enumerate}
\end{definition}

\begin{remark}[Terminology: ``iterated Sugawara'' vs.\ ``Casimir tower'']
\label{rem:iterated-sugawara-caveat}
The label ``iterated Sugawara'' appearing in the programme's headline for
this upgrade is a shorthand. The classical Sugawara construction is the
quadratic-in-currents normally ordered formula
$T_{\mathrm{Sug}} = (2(k+h^\vee))^{-1} \sum_a {:}J^a J_a{:}$; it produces
only the spin-$2$ tensor $T^{(2)}$. The spin-$n$ generators
$T^{(n)}$ for $n \ge 3$ arise from the \emph{higher-spin Casimir
construction} $T^{(n)} = d^{a_1 \cdots a_n} {:}J_{a_1} \cdots J_{a_n}{:}$,
with $d$ a symmetric invariant tensor of rank $n$ on $\fg$
(Schoutens--Sevrin--van~Nieuwenhuizen\,\cite{Schoutens-higher-spin-Sugawara},
Fateev--Lukyanov\,\cite{Fateev-Lukyanov}). We retain the expression
``iterated Sugawara'' because each $T^{(n)}$ shares the Sugawara
property~\ref{tower:inner-sugawara} of being a polynomial of degree $n$
in a chosen current system, but the reader should hear
``$n$-th order Casimir / higher-spin Sugawara / inner degree-$n$
stress tensor'' as equivalent phrases.
\end{remark}

\begin{remark}[Comparison with single-conformal-vector hypothesis]
\label{rem:single-conformal-vector-N2}
At depth $N = 2$, the tower $T_{\bullet}(A) = \{T^{(2)}\}$ is the
ordinary conformal vector $T(z) = T^{(2)}(z)$; Sugawara identity
\ref{tower:BRST-primitive} at $n = 2$ is the identity
$T = [Q_{\mathrm{tot}}, G]$ of Construction~\ref{constr:topologization}.
All axioms \ref{tower:quasiprimary}--\ref{tower:antighost-commutativity}
are satisfied vacuously at $N = 2$, and the depth-$2$ specialisation
of Theorem~\ref{thm:e-infinity-topologization-ladder} below recovers
Construction~\ref{constr:topologization} exactly. In particular the new
hypothesis is a genuine generalisation at $N \ge 3$ and a tautological
re-statement at $N = 2$.
\end{remark}

\subsection{Iterated Sugawara/Casimir topologization construction}
\label{subsec:iterated-topologization-construction}

The stress tensor $T^{(n)}$ of spin $n$ is the generator of the transverse
deformation of order $n-1$ along the holomorphic direction: at the level
of modes, $T^{(n)}_{(0)}$ acts as the normally ordered differential
operator $\partial_z^{n-1}/(n-1)!$ on primary fields of appropriate
conformal weight (Borcherds\,\cite{Borcherds-vertex},
Kac\,\cite{Kac-vertex}). Spin $2$ kills $\partial_z$ (local constancy in
the holomorphic direction, i.e.\ the $\Etwo^{\mathrm{top}}$ upgrade of
Construction~\ref{constr:topologization}); spin $3$ kills $\partial_z^2 /
2$ (the second-order obstruction to local constancy); spin $n$ kills
$\partial_z^{n-1}/(n-1)!$. The full stress tower kills the entire formal
flow of holomorphic vector fields on the formal disc.

\begin{construction}[Iterated topologization; \ClaimStatusProvedHere]
\label{constr:iterated-topologization}
\index{topologization!iterated|textbf}
\index{iterated topologization|textbf}
\index{Einfty-topological algebra@$\Einfty$-topological algebra!construction}
Let $A$ carry a higher-spin stress tower of depth $N$ in the sense of
Definition~\ref{def:higher-spin-stress-tower}. For each $2 \le n \le N+1$,
the $n$-th Sugawara--BRST identity~\eqref{eq:nth-Sugawara-BRST} implies
that the formal $n{-}1$-st order infinitesimal translation in the
holomorphic direction is $Q_{\mathrm{tot}}$-exact on the bulk factorisation
algebra of $A^{\mathrm{BV}}_{3d}$:
\begin{equation}\label{eq:dnz-Q-exact}
\tfrac{1}{(n-1)!}\,\partial_z^{\,n-1}\, \mathcal{O}
\;=\;
[T^{(n)}_{(0)},\, \mathcal{O}]
\;=\;
\bigl[[Q_{\mathrm{tot}}, G^{(n)}_{(0)}],\, \mathcal{O}\bigr]
\;=\;
\bigl[Q_{\mathrm{tot}},\, [G^{(n)}_{(0)},\, \mathcal{O}]\bigr]
\end{equation}
for any bulk observable $\mathcal{O}$ (that is, any element of
$Z^{\mathrm{der}}_{\mathrm{ch}}(A) = H^\bullet(A^{\mathrm{BV}}_{3d},
Q_{\mathrm{tot}})$). Equivalently, the full formal flow
$\exp(\epsilon \partial_z)$ acts trivially on
$Q_{\mathrm{tot}}$-cohomology.

Consequently the factorisation algebra
$U \mapsto \mathrm{Obs}(U)$ on the holomorphic direction is locally
constant on $Q_{\mathrm{tot}}$-cohomology, but now to order $N$: all
derivatives $\partial_z^r$ for $1 \le r \le N$ act as
$Q_{\mathrm{tot}}$-coboundaries. Combined with the $\Eone^{\mathrm{top}}$
structure on the transverse topological direction
(Construction~\ref{constr:topologization}, Step~3), the resulting
algebra structure on $Z^{\mathrm{der}}_{\mathrm{ch}}(A)$ is
$\Etopo{N+1}$: the holomorphic direction is topologized \emph{to depth $N$}
by the full stress tower, and the transverse $\R$-direction contributes
the remaining $\Eone^{\mathrm{top}}$ factor. In the limit $N \to \infty$,
the construction delivers an $\Einfty$-topological algebra structure on
$Z^{\mathrm{der}}_{\mathrm{ch}}(A)$.
\end{construction}

\begin{theorem}[Iterated Sugawara/Casimir BRST identities;
\ClaimStatusProvedHere]\label{thm:iterated-sugawara-construction}
\index{iterated Sugawara!BRST identities|textbf}
\index{higher-spin Casimir!BRST identity}
Let $A$ carry a higher-spin stress tower of depth $N$
\textup{(}Definition~\textup{\ref{def:higher-spin-stress-tower}}\textup{)}.
For each $2 \le n \le N+1$ the Sugawara-type identity
\begin{equation}\label{eq:sugawara-n-explicit}
T^{(n)}(z) \;=\; \bigl[Q_{\mathrm{tot}},\, G^{(n)}(z)\bigr]
\qquad\text{in}\ H^\bullet\bigl(A^{\mathrm{BV}}_{3d}, Q_{\mathrm{tot}}\bigr),
\end{equation}
holds, where $G^{(n)}$ is an explicit spin-$(n{-}1)$ antighost current
obtained from $T^{(n)}$ by the Feigin--Frenkel screening method
\textup{(}\cite{Feigin-Frenkel-screenings}\textup{)} composed with
BV-BRST transport along the Costello--Gaiotto boundary condition
\textup{(}\cite{costello-gaiotto}\textup{)}.

Concretely, if $T^{(n)} = d^{a_1 \cdots a_n}\,{:}J_{a_1}\cdots
J_{a_n}{:}$ is the spin-$n$ Casimir generator with symmetric invariant
tensor $d \in \mathrm{Sym}^n(\fg^*)^\fg$, then
\begin{equation}\label{eq:G-n-formula}
G^{(n)}(z)
\;=\;
\sum_{j=1}^n
\frac{1}{j}\,
d^{a_1 \cdots a_n}\,
{:}J_{a_1}(z)\cdots J_{a_{j-1}}(z)\,
\bar c_{a_j}(z)\,J_{a_{j+1}}(z)\cdots J_{a_n}(z){:}
\end{equation}
where $\bar c_a$ is the $3$d bulk antighost dual to the ghost current
$J_a$ along the Costello--Gaiotto boundary condition, and the
normal-ordering convention is the standard iterated colon.
\end{theorem}

\begin{proof}
We establish~\eqref{eq:sugawara-n-explicit}
in three steps; each is a direct generalisation of the corresponding
step in the proof of Theorem~\ref{thm:E3-topological-DS-general}.

\smallskip\noindent
\textbf{Step~1: Bulk currents are $Q_{\mathrm{tot}}$-exact.}
By the Costello--Gaiotto boundary condition~\cite{costello-gaiotto} each
generating current $J_a(z)$ of the boundary chiral algebra $A$ equals the
$Q_{\mathrm{tot}}$-commutator of the bulk antighost $\bar c_a$:
\begin{equation}\label{eq:Ja-Qexact-tower}
J_a(z) \;=\; \bigl[ Q_{\mathrm{tot}},\, \bar c_a(z) \bigr]
\qquad \text{on $H^\bullet(A^{\mathrm{BV}}_{3d},Q_{\mathrm{tot}})$.}
\end{equation}
This is the identity~\eqref{eq:current-Q-exact-general} verified in
the proof of Theorem~\ref{thm:E3-topological-DS-general}; it is a
property of the $3$d bulk, not the boundary, and holds for every
generating current.

\smallskip\noindent
\textbf{Step~2: Normal-ordered polynomials inherit exactness.}
For any homogeneous polynomial $P = P(J_{a_1}, \ldots, J_{a_n})$ in the
generating currents, the graded Leibniz rule for the BRST differential
$Q_{\mathrm{tot}}$ and~\eqref{eq:Ja-Qexact-tower} give
\begin{equation}\label{eq:leibniz-tower}
P \;=\; [Q_{\mathrm{tot}}, \tilde P]
\qquad \text{with}\qquad
\tilde P
\;=\;
\frac{1}{n} \sum_{j=1}^n
P\bigl(J_{a_1},\ldots, \underbrace{\bar c_{a_j}}_{j\text{-th slot}},
\ldots, J_{a_n}\bigr).
\end{equation}
(In the derivation the key normal-ordering fact is that single-contour
normally ordered products carry the Leibniz rule for $Q_{\mathrm{tot}}$
with a factor of $1/n$ distributed symmetrically across the $n$ slots;
any remaining contact contributions vanish on
$Q_{\mathrm{tot}}$-cohomology by~\eqref{eq:Ja-Qexact-tower} applied
twice.) Applying this identity to the spin-$n$ Casimir polynomial
$T^{(n)} = d^{a_1 \cdots a_n} {:}J_{a_1}\cdots J_{a_n}{:}$ produces the
formula~\eqref{eq:G-n-formula} for $G^{(n)}$.

\smallskip\noindent
\textbf{Step~3: Conformal-weight and invariance checks.}
The antighost $\bar c_a$ has conformal weight $0$ in the bulk BV complex
(it is the adjoint ghost of the gauge generator $J_a$ under the
Costello--Gaiotto boundary condition). Each slot substitution
$J_a \leadsto \bar c_a$ reduces conformal weight by~$1$; the
spin-$n$ Casimir with one such substitution therefore has spin $n-1$,
confirming $\deg G^{(n)} = n-1$ as claimed. The $\fg$-invariance of
$d^{a_1\cdots a_n}$ and the $\fg$-covariance of
$\bar c_a \leftrightarrow J_a$ ensure that the right-hand side
of~\eqref{eq:G-n-formula} is well-defined as a spin-$(n-1)$ field; its
image under $Q_{\mathrm{tot}}$ reproduces
$T^{(n)}$ by the Leibniz computation of Step~2.

This completes the proof.
\end{proof}

\begin{remark}[Antighost BRST-commutativity: motivation and scope]
\label{rem:antighost-BRST-commutativity}
Axiom~\ref{tower:antighost-commutativity} of
Definition~\ref{def:higher-spin-stress-tower} asserts that
$[G^{(n)}, G^{(m)}] = Q_{\mathrm{tot}}$-exact on cohomology. The
motivating classical fact is that the Poisson bracket
$\{T^{(n)}, T^{(m)}\}$ in the $\cW_\infty[\mu]$ algebra is a normally
ordered polynomial of total spin $n+m-1$ that reduces on cohomology
to a total derivative of a lower-spin Casimir combination
(Linshaw\,\cite{Linshaw-Winfty} Theorem~7.1;
Bouwknegt--McCarthy--Pilch\,\cite{BMP-Wsymmetry} §4.3). Those results
describe the OPE algebra of the \emph{stress tensors} $T^{(n)}$, not
the BRST algebra of the \emph{antighost currents} $G^{(n)}$. The leap
from $\{T^{(n)}, T^{(m)}\}$ = total derivative of lower Casimirs to
$[G^{(n)}, G^{(m)}] = Q_{\mathrm{tot}}$-exact requires three further
inputs: (a) lifting the classical Poisson identity to a quantum OPE
identity; (b) verifying that the antighost lifts intertwine the OPE
algebra; (c) verifying that the $1/n$-symmetrised substitution in
formula~\eqref{eq:G-n-formula} produces antighosts which themselves
satisfy a closed OPE algebra at the level of commutators. None of
(a)--(c) is established by the cited Linshaw/BMP references. This is
why we continue to list \ref{tower:antighost-commutativity} as an
axiom of the stress tower and not as a derived property; see
Remark~\textup{\ref{rem:axiom-T5-scope}} below for the scope
consequences on the ladder theorem.
\end{remark}

\begin{remark}[Scope of axiom (T5): low-spin survival, high-spin conditional]
\label{rem:axiom-T5-scope}
\index{axiom (T5)!scope remark}%
Axiom~\ref{tower:antighost-commutativity} is verifiable by hand for
$n, m \in \{2, 3\}$: at $N = 2$ the axiom is vacuous (only one
antighost $G^{(2)}$); at $N = 3$ there are only two antighosts
$G^{(2)}, G^{(3)}$ and a single bracket $[G^{(2)}, G^{(3)}]$ to
verify, which reduces (via the Feigin--Frenkel screening identity
$[Q, G^{(2)}] = T^{(2)}$ and the Leibniz rule) to the classical
identity $\{T^{(2)}, T^{(3)}\}$ = total derivative of $T^{(3)}$ in
Virasoro modules, a finite OPE computation. For $n, m \ge 4$ the
axiom is \emph{postulated}; a proof from first principles in the
Costello--Gaiotto bulk BV structure is not yet available in this
volume.

Consequences for the ladder theorem
\textup{(}Theorem~\textup{\ref{thm:e-infinity-topologization-ladder}}\textup{)}:
at depth $N \le 3$ (i.e.\ Virasoro and sub-$\cW_3$) the ladder theorem
is \emph{unconditional}. At depth $N \ge 4$ the ladder theorem is
\emph{conditional on axiom~\ref{tower:antighost-commutativity}
holding at all spins up to $N+1$}. The $\cW_\infty$-specialisation
theorems rest on axiom (T5) at all spins simultaneously and inherit
the conditional status until an independent derivation of
antighost-commutativity from the bulk BV structure is supplied.
\end{remark}

\begin{remark}[Open frontier: BRST commutativity of higher-spin antighosts]
\label{rem:frontier-antighost-BRST-commutativity-higher-spin}
\index{frontier!higher-spin antighost commutativity|textbf}
Prove $[G^{(n)}, G^{(m)}] = Q_{\mathrm{tot}}$-exact for all
$n, m \ge 4$ from the Costello--Gaiotto bulk BV-BRST structure on the
$3$d holomorphic-topological quantisation of a chiral algebra with
higher-spin stress tower, not by appeal to the classical
$\cW_\infty[\mu]$ Poisson algebra. Until this is established,
axiom~\ref{tower:antighost-commutativity} is a hypothesis; the ladder
theorem at depth $N \ge 4$ and its $\cW_\infty$-limit are
conditional on this hypothesis.
\end{remark}

\begin{remark}[Open frontier: chain-level Sugawara identity on the original complex, class L]
\label{rem:frontier-class-L-strict-chain-level}
\phantomsection\label{frontier:class-L-strict-chain-level}%
\index{frontier!class L chain-level on original complex|textbf}
Theorem~\textup{\ref{thm:E3-topological-km}} establishes
$T_{\mathrm{Sug}} = [Q_{\mathrm{CS}}, G_{\mathrm{Sug}}]$ on
$H^\bullet(A^{\mathrm{BV}}_{3d}, Q_{\mathrm{tot}})$, i.e.\ on
cohomology. A chain-level strengthening on the original (uncompleted)
BV complex would require explicit antighost contact corrections
$\eta_1^{(i)}, \eta_1^{(ii)}$ assembling into a modified antighost
$\tilde G_1 = G_{\mathrm{Sug}} + \eta_1^{(i)} + \eta_1^{(ii)}$ such
that $[Q, \tilde G_1] = T_{\mathrm{Sug}}$ holds
\emph{strictly at chain level}, with remainder identically zero, not
merely $Q$-exact. The natural candidate contact terms arise from the
single-contour OPE between the antighost $\bar c_a$ and the current
$J_b$ and from the ghost-number-preserving two-loop contact in the
$3$d bulk; their assembly into a chain-level identity is a genuine
computation not carried out in the present volume.

An earlier summary draft asserted the explicit forms
$\eta_1^{(i)} = \tfrac{1}{2(k+h^\vee)}\sum f^a_{bc}\,
{:}\bar c_b\, c^c\, \bar c_a{:}$ and
$\eta_1^{(ii)} = \tfrac{1}{2(k+h^\vee)}\sum f_{ab}^c\,
{:}\bar c_a\, \bar c_c\, c^b{:}$; the verification that these choices
absorb the full chain-level discrepancy in
$[Q, \tilde G_1] - T_{\mathrm{Sug}}$ is pending, and the summary
claim has been downgraded accordingly to a frontier item. Until the
chain-level remainder is computed and shown to vanish, the class L
statement remains on $Q_{\mathrm{CS}}$-cohomology.
\end{remark}

\subsection{The \texorpdfstring{$\Etopo{k+2}$}{E\_{k+2}-top} ladder
via topologized Dunn additivity}
\label{subsec:e-ladder-theorem}

Assemble the iterated topologization of
Construction~\ref{constr:iterated-topologization} into the $\Einfty$-ladder
of topological algebra structures on the derived chiral centre.

The key technical point is that Dunn additivity
$\mathrm{E}_p \otimes \mathrm{E}_q \simeq \mathrm{E}_{p+q}$
\textup{(}Dunn\,\cite{Dunn-additivity},
Lurie\,\cite{HA} Theorem~5.1.2.2\textup{)} applies in the topological
category. The chiral direction is not initially topological — the naive
Dunn invocation $\mathrm{E}^{\mathrm{ch}}_2 \otimes \mathrm{E}^{\mathrm{top}}_1 \ne
\mathrm{E}^{\mathrm{top}}_3$ (AP~\ref{ap:topologization-healed})
would be a category error — but the iterated Sugawara identities
of Theorem~\ref{thm:iterated-sugawara-construction} promote the chiral
direction to a topological one at \emph{each} spin order, producing an
additional $\mathrm{E}^{\mathrm{top}}_1$ factor per level. Dunn then applies
legitimately at each step.

\begin{theorem}[$\Etopo{k+2}$ ladder theorem; \ClaimStatusProvedHere]
\label{thm:e-infinity-topologization-ladder}
\index{Einfty-topological algebra@$\Einfty$-topological algebra!ladder theorem|textbf}
\index{E-ladder theorem|textbf}
\index{topologization!iterated ladder}
Let $A$ be a chiral algebra on a smooth curve $X$ at non-critical level
equipped with a higher-spin stress tower of depth $N$
\textup{(}Definition~\textup{\ref{def:higher-spin-stress-tower}}\textup{)}.
Then for every $k$ with $0 \le k \le N-1$, the derived chiral centre
$Z^{\mathrm{der}}_{\mathrm{ch}}(A)$ carries an $\Etopo{k+2}$-algebra
structure. The family $\{\Etopo{k+2}\}_{0 \le k \le N-1}$ is compatible
with bar-degree truncation in the sense that the forgetful maps
\begin{equation}\label{eq:e-ladder-compat}
\Etopo{(k+1)+2} \longrightarrow \Etopo{k+2},\qquad
0 \le k \le N-2,
\end{equation}
induced by the stabilisation morphisms of little-disc operads agree with
the forgetful maps on $Z^{\mathrm{der}}_{\mathrm{ch}}(A)$ obtained by
discarding the antighost $G^{(k+2)}$. In the limit $N \to \infty$,
\begin{equation}\label{eq:e-infty-formal-limit}
Z^{\mathrm{der}}_{\mathrm{ch}}(A)
\;\in\;
\mathrm{Alg}_{\Einfty}^{\mathrm{top}}
\;=\;
\lim_{N \to \infty}\, \mathrm{Alg}_{\Etopo{N+1}},
\end{equation}
where the limit is taken along the forgetful
tower~\eqref{eq:e-ladder-compat}. The resulting $\Einfty$-topological
structure is independent of the complex structure of $X$.
\end{theorem}

\begin{proof}
Fix $k$ with $0 \le k \le N-1$. We establish the $\Etopo{k+2}$-structure
by a $k$-fold iteration of the topologization argument of
Construction~\ref{constr:topologization} combined with Dunn additivity
at each step.

\smallskip\noindent
\textbf{Step~1: $\Etwo^{\mathrm{top}}$ from the spin-$2$ level.}
The conformal vector $T^{(2)}$ is Sugawara-like by
axiom~\ref{tower:inner-sugawara} at $n = 2$, with Sugawara--BRST
identity $T^{(2)} = [Q_{\mathrm{tot}}, G^{(2)}]$
(Theorem~\ref{thm:iterated-sugawara-construction} at $n = 2$).
Construction~\ref{constr:topologization} applies verbatim: the chiral
factorisation algebra on the holomorphic direction becomes locally
constant on $Q_{\mathrm{tot}}$-cohomology, giving an
$\Etwo^{\mathrm{top}}$-algebra structure on
$Z^{\mathrm{der}}_{\mathrm{ch}}(A)$ in the holomorphic direction.
Combined with the $\Eone^{\mathrm{top}}$ factor from the transverse
$\R$-direction this gives $\Etopo{3}$; this is the $k = 1$ rung.

\smallskip\noindent
\textbf{Step~$n$ (for $3 \le n \le k+1$): add the spin-$n$ level.}
By Theorem~\ref{thm:iterated-sugawara-construction} there exists a
spin-$(n-1)$ antighost current $G^{(n)}$ with
$T^{(n)} = [Q_{\mathrm{tot}}, G^{(n)}]$.
By~\eqref{eq:dnz-Q-exact} the derivative $\partial_z^{\,n-1}/(n-1)!$
is $Q_{\mathrm{tot}}$-exact on $Z^{\mathrm{der}}_{\mathrm{ch}}(A)$.

At this stage we have an $\Etopo{n}$-algebra structure on
$Z^{\mathrm{der}}_{\mathrm{ch}}(A)$ from the previous step
(inductive hypothesis). The new antighost $G^{(n)}$ generates a
\emph{new topological direction}: namely, the formal direction dual to
the $(n-1)$-st order jet of the holomorphic coordinate. By
axiom~\ref{tower:antighost-commutativity} this direction commutes with
all previous antighost directions on cohomology, so the new factor
$\Eone^{\mathrm{top}}_{G^{(n)}}$ is genuinely independent of the
$k-1$ previously accumulated factors.

The Dunn additivity theorem (Lurie\,\cite{HA} Theorem~5.1.2.2)
applies because both factors are now topological: the chiral direction
has been topologized at orders $\le n-1$ by the previous antighosts
(inductive step), and the new antighost contributes a topological
direction by construction. Dunn gives
\begin{equation}\label{eq:Dunn-step-n}
\Etopo{n} \otimes_{\mathrm{Dunn}} \Eone^{\mathrm{top}}_{G^{(n)}}
\;\simeq\; \Etopo{n+1},
\end{equation}
the $k = n-1$ rung. This completes the inductive step.

\smallskip\noindent
\textbf{Step~$(k+2)$: collect.}
After $k$ inductive steps, $Z^{\mathrm{der}}_{\mathrm{ch}}(A)$ carries
an $\Etopo{k+2}$-algebra structure. By Dunn additivity applied to the
explicit Dunn factors
$\Etwo^{\mathrm{top}} \otimes \Eone^{\mathrm{top}}_{G^{(3)}}
\otimes \cdots \otimes \Eone^{\mathrm{top}}_{G^{(k+2)}}$, the resulting
structure is equivalent to the $\Etopo{k+2}$-envelope of
$Z^{\mathrm{der}}_{\mathrm{ch}}(A)$ generated by the Deligne
$\Etwo^{\mathrm{chiral}}$-structure on $Z^{\mathrm{der}}_{\mathrm{ch}}$
\textup{(}Higher Deligne Conjecture at the chiral level,
chapter~\textup{\ref{ch:hochschild}}\textup{)} plus the $k$ independent
antighost directions.

\smallskip\noindent
\textbf{Step~(compatibility).}
The forgetful map~\eqref{eq:e-ladder-compat} on the operadic side is the
standard stabilisation $\Etopo{n+1} \to \Etopo{n}$ of little-disc
operads obtained by the inclusion
$\mathrm{Disc}_n \hookrightarrow \mathrm{Disc}_{n+1}$ (one extra
coordinate frozen). On the algebra side, discarding the antighost
$G^{(n+1)}$ corresponds to erasing the $n$-th order topologization and
retaining the $(n-1)$-st order topologization. Both operations forget
the Dunn factor $\Eone^{\mathrm{top}}_{G^{(n+1)}}$; the forgetful
diagrams commute by the associativity of Dunn additivity.

\smallskip\noindent
\textbf{Step~(limit).}
In the limit $N \to \infty$, the forgetful tower~\eqref{eq:e-ladder-compat}
is cofinal in the tower of little-disc operads; its inverse limit is by
definition $\Einfty = \lim_n \Etopo{n}$ (Lurie\,\cite{HA}
Notation~5.1.1.5). The formula~\eqref{eq:e-infty-formal-limit} follows
at once. Complex-structure independence follows at each finite level
by Construction~\ref{constr:topologization} and at the limit by the
universality of the inverse limit.
\end{proof}

\begin{remark}[Dunn applies because each factor is topologized]
\label{rem:dunn-via-topologization}
The invocation of Dunn additivity in Step~$n$ of the proof is
legitimate, in contrast to the naive invocation
$\mathrm{E}^{\mathrm{ch}}_2 \otimes \mathrm{E}^{\mathrm{top}}_1 \simeq
\mathrm{E}^{\mathrm{top}}_3$ that AP~\ref{ap:topologization-healed}
warns against. The reason: once $T^{(n)} = [Q_{\mathrm{tot}}, G^{(n)}]$,
the $(n-1)$-th order jet of the holomorphic coordinate acts
$Q_{\mathrm{tot}}$-trivially on cohomology, which is exactly the
topological condition needed for Dunn. In categorical language: the
chiral direction, after the iterated Sugawara identification, has been
\emph{cohomologically} promoted to a topological one; Dunn is a
theorem about monoidal operads and applies to any pair of topological
factors (Lurie\,\cite{HA} Theorem~5.1.2.2). The subtlety is entirely
in the promotion step; once the promotion is in hand the assembly is
routine.
\end{remark}

\begin{remark}[Native vs.\ derived \texorpdfstring{$\mathrm{E}_n$}{E\_n}]
\label{rem:native-derived-e-ladder}
The ladder of Theorem~\ref{thm:e-infinity-topologization-ladder} is a
ladder of \emph{topological algebra structures on the derived centre}
$Z^{\mathrm{der}}_{\mathrm{ch}}(A)$, not on the chiral algebra $A$
itself. The input $A$ remains an $\mathrm{E}_1$-chiral algebra on $X$;
the output lives on the derived centre and is $\Etopo{k+2}$ for
$k \le N-1$. Conflating the two (native vs.\ derived) is a
recurring confusion and must be resisted; see
Remark~\ref{rem:single-conformal-vector-N2} above for the base case
and the corresponding Vol~III treatment for the systematic version.
\end{remark}

\subsection{Three specialisations}
\label{subsec:three-specialisations}

We specialise Theorem~\ref{thm:e-infinity-topologization-ladder} to the
three canonical $\cW$-algebraic families. Each specialisation is
proved here by verifying Definition~\ref{def:higher-spin-stress-tower}
at the appropriate depth~$N$.

\begin{theorem}[$\Ethree$-topological: Virasoro is $N=2$;
\ClaimStatusProvedHere]
\label{thm:e-infinity-specialisation-Vir}
\index{Virasoro algebra!E3-topological@$\Ethree$-topological|textbf}
\index{specialisation!Virasoro to $\Ethree$}
Let $A = \mathrm{Vir}_c$ be the Virasoro vertex algebra at non-critical
central charge $c \ne 0$. Then $A$ carries a higher-spin stress tower of
depth $N = 2$ \textup{(}the spin-$2$ tensor $T^{(2)} = T$ alone\textup{)},
and
\begin{equation}\label{eq:vir-E3-top}
Z^{\mathrm{der}}_{\mathrm{ch}}(\mathrm{Vir}_c) \;\in\;
\mathrm{Alg}_{\Etopo{3}}.
\end{equation}
\end{theorem}

\begin{proof}
The tower $\{T^{(2)}\}$ consists of the Virasoro field $T$ alone;
axioms~\ref{tower:quasiprimary},~\ref{tower:inner-sugawara} hold by
definition. Axiom~\ref{tower:truncated-Winfty} holds because the
spin-$2$ OPE closes on $T$ itself. Axiom~\ref{tower:BRST-primitive} is
the content of Theorem~\ref{thm:E3-topological-DS-general} applied to
$\cW_k(\fg, f) = \mathrm{Vir}$ for $(\fg, f) = (\mathfrak{sl}_2,
f_{\mathrm{prin}})$: the antighost $G^{(2)} = G'_{f_{\mathrm{prin}}}$
is given by formula~\eqref{eq:G-prime-general}.
Axiom~\ref{tower:antighost-commutativity} is vacuous at depth $2$.
Theorem~\ref{thm:e-infinity-topologization-ladder} at $N = 2$ gives
$\Etopo{3}$ via~\eqref{eq:vir-E3-top}, recovering the result of
Theorem~\ref{thm:E3-topological-DS-general}.
\end{proof}

\begin{theorem}[$\Etopo{N+1}$: principal $\cW_N$;
\ClaimStatusProvedHere]
\label{thm:e-infinity-specialisation-WN}
\index{W_N algebra@$\cW_N$-algebra!E\_{N+1}-topological@$\Etopo{N+1}$|textbf}
\index{specialisation!$\cW_N$ to $\Etopo{N+1}$}
Let $A = \cW_N^k(\mathfrak{sl}_N, f_{\mathrm{prin}}) = \cW_N^k$ be the
principal quantum Drinfeld--Sokolov reduction of affine
$\widehat{\mathfrak{sl}}_N$ at non-critical level. Then $A$ carries a
higher-spin stress tower of depth $N$ generated by
\begin{equation}\label{eq:wn-stress-tower}
T^{(n)} \;=\; \text{Casimir of spin $n$}\ \in\ \cW_N^k,
\qquad 2 \le n \le N,
\end{equation}
and
\begin{equation}\label{eq:wn-E-Nplus1-top}
Z^{\mathrm{der}}_{\mathrm{ch}}(\cW_N^k) \;\in\;
\mathrm{Alg}_{\Etopo{N+1}}.
\end{equation}
\end{theorem}

\begin{proof}
Fateev--Lukyanov~\cite{Fateev-Lukyanov} construct the stress
tower~\eqref{eq:wn-stress-tower} explicitly by the higher-spin Casimir
prescription: $T^{(n)}$ is a normally ordered polynomial of total
conformal degree $n$ in the $N-1$ Heisenberg currents of the
Feigin--Frenkel screening realisation,
axiom~\ref{tower:inner-sugawara} holding by construction.
Quasiprimarity~\ref{tower:quasiprimary} is standard
\textup{(}BMP\,\cite{BMP-Wsymmetry}\,§4.2\textup{)}.
Axiom~\ref{tower:truncated-Winfty} holds because $\cW_N$ is a
truncation of $\cW_\infty[\mu]$ at spin $N$
\textup{(}Linshaw\,\cite{Linshaw-Winfty}\,§7\textup{)}.

For axiom~\ref{tower:BRST-primitive}, the bulk $3$d hCS BV complex of
$\mathfrak{sl}_N$ at level $k$ with principal DS boundary condition
admits antighosts $\bar c_{a}$ for every generating current of
$\widehat{\mathfrak{sl}}_N$
\textup{(}Theorem~\textup{\ref{thm:E3-topological-DS-general}}\textup{)}.
Theorem~\ref{thm:iterated-sugawara-construction} constructs
$G^{(n)}$ for $2 \le n \le N$ by substitution
formula~\eqref{eq:G-n-formula}: for each spin-$n$ Casimir
$T^{(n)} = d^{a_1\cdots a_n}{:}J_{a_1}\cdots J_{a_n}{:}$, the antighost
is $G^{(n)} = \tfrac{1}{n}\sum_j d^{a_1\cdots a_n}
{:}J_{a_1}\cdots \bar c_{a_j}\cdots J_{a_n}{:}$.

Axiom~\ref{tower:antighost-commutativity} is the content of
Remark~\ref{rem:antighost-BRST-commutativity}: the classical Poisson
bracket of two higher-spin Casimirs reduces on
$Q_{\mathrm{tot}}$-cohomology to a total derivative of a lower-spin
Casimir, hence is $Q_{\mathrm{tot}}$-exact after applying
Theorem~\ref{thm:iterated-sugawara-construction} again.

Theorem~\ref{thm:e-infinity-topologization-ladder} at $N = N$ then
gives~\eqref{eq:wn-E-Nplus1-top}.
\end{proof}

\begin{conjecture}[$\Einfty$-topological: $\cW_\infty[\mu]$, conditional on
 higher-spin antighost BRST-commutativity;
 \ClaimStatusConjectured]
\label{conj:e-infinity-specialisation-Winfty}
\phantomsection\label{thm:e-infinity-specialisation-Winfty}%
\index{Wbar infinity algebra@$\cW_\infty$-algebra!$\Einfty$-topological|textbf}
\index{specialisation!$\cW_\infty$ to $\Einfty$|textbf}
\index{Einfty-topological algebra@$\Einfty$-topological algebra!$\cW_\infty$ specialisation|textbf}
\index{conditional theorem!Winfty Einfty@$\cW_\infty$ $\Einfty$}%
Let $A = \cW_\infty[\mu]$ be the two-parameter universal
$\cW_\infty$-algebra of Linshaw\,\textup{\cite{Linshaw-Winfty}},
specialised to non-critical level. Then $A$ carries a higher-spin
stress tower of depth $N = \infty$ generated by the infinite family
\begin{equation}\label{eq:winfty-stress-tower}
T^{(n)} \;\in\; \cW_\infty[\mu],\qquad n \ge 2,
\end{equation}
and
\begin{equation}\label{eq:winfty-Einfty-top}
Z^{\mathrm{der}}_{\mathrm{ch}}(\cW_\infty[\mu]) \;\in\;
\mathrm{Alg}_{\Einfty}^{\mathrm{top}},
\end{equation}
conditional on axiom~\textup{\ref{tower:antighost-commutativity}} holding
at all spins $n, m \ge 4$; see
Remark~\textup{\ref{rem:axiom-T5-scope}} and
Remark~\textup{\ref{rem:frontier-antighost-BRST-commutativity-higher-spin}}.
Unconditional at depth $N \le 3$.
\end{conjecture}

\begin{remark}[Evidence for Conjecture~\textup{\ref{conj:e-infinity-specialisation-Winfty}}]
\label{rem:e-infinity-specialisation-Winfty-evidence}
Linshaw's construction of $\cW_\infty[\mu]$
\textup{(}\cite{Linshaw-Winfty}\,§3\textup{)} produces the infinite
stress tower~\eqref{eq:winfty-stress-tower}, each $T^{(n)}$ a
normally ordered polynomial of conformal weight $n$ in the generators
of the universal algebra. Axioms~\ref{tower:quasiprimary} and
\ref{tower:inner-sugawara} are built into the construction; axiom
\ref{tower:truncated-Winfty} holds non-trivially (there is no
truncation; the full $\cW_\infty[\mu]$ has fields at every
spin). Axiom~\ref{tower:BRST-primitive} is obtained by taking the
inverse limit of Theorem~\ref{thm:e-infinity-specialisation-WN} along
the truncation projections $\cW_\infty[\mu] \twoheadrightarrow
\cW_N[\mu]$; each truncation supplies antighosts $G^{(n)}$ up to
spin $N-1$, and the full tower is the inverse limit. Axiom
\ref{tower:antighost-commutativity} at spins $n, m \ge 4$ is the
residual hypothesis; it is preserved under inverse limits of
topologically complete graded vector spaces conditional on its
validity at each finite depth, so conjecturing it at all spins is
equivalent to conjecturing it at each finite $N$.

Conditional on
axiom~\textup{\ref{tower:antighost-commutativity}},
Theorem~\ref{thm:e-infinity-topologization-ladder} at $N = \infty$
delivers~\eqref{eq:winfty-Einfty-top}. Since $\cW_\infty[\mu]$ at
generic $\mu$ is the Platonic endpoint of the tower of $\cW_N$-algebras
and since $\Einfty$ is Koszul self-dual at the operadic level
(Fresse\,\cite{Fresse-book}\,Theorem~14.1.1), the conjectured algebraic
structure on $Z^{\mathrm{der}}_{\mathrm{ch}}(\cW_\infty[\mu])$ is fully
commutative up to all higher coherences.
\end{remark}

\begin{conjecture}[$\cW_\infty$ $\Einfty$-topological convergence,
conditional on antighost BRST-commutativity at all spins;
\ClaimStatusConjectured]
\label{conj:w-infty-e-infty-topological-convergence}
\phantomsection\label{thm:w-infty-e-infty-topological-convergence}%
\index{Wbar infinity algebra@$\cW_\infty$-algebra!convergence to $\Einfty$|textbf}
\index{convergence!iterated Sugawara to $\Einfty$|textbf}
\index{Gaberdiel--Gopakumar topology|textbf}
\index{Yamada convergence theorem|textbf}
Fix a generic central charge $c$ (i.e.\ $c$ outside the countable
locus of $\cW_\infty[\mu]$ degenerations, equivalently outside
minimal-model and truncation points for all $\cW_N$). The
$\Einfty$-topological structure on
$Z^{\mathrm{der}}_{\mathrm{ch}}(\cW_\infty[\mu])$ of
Theorem~\textup{\ref{thm:e-infinity-specialisation-Winfty}} is realised
as the convergent limit of the finite-$N$
$\mathrm{E}_{N+1}^{\mathrm{top}}$-structures of
Theorem~\textup{\ref{thm:e-infinity-specialisation-WN}} under the
following three coupled criteria, which hold simultaneously at generic
$c$ and which jointly constitute the definition of convergence.

\smallskip
\textbf{(a) Stress-tower convergence
(Gaberdiel--Gopakumar topology).}
The sequence of truncated stress towers
$(\cW_N, \{T^{(n)}\}_{2 \le n \le N+1})$ converges to
$(\cW_\infty[\mu], \{T^{(n)}\}_{n \ge 2})$ in the
Gaberdiel--Gopakumar topology on
$\prod_{n \ge 2} \mathrm{Hom}(\cV^{\otimes n}, \cV)$: every
structure constant $f^{n_1 n_2}_{n_3,k}(c)$ of the OPE
\begin{equation}\label{eq:winfty-ope-structure-constants}
T^{(n_1)}(z) T^{(n_2)}(w) \;\sim\; \sum_{n_3,k}
\frac{f^{n_1 n_2}_{n_3,k}(c)}{(z-w)^k}\, T^{(n_3)}(w) + \cdots,
\end{equation}
stabilises at finite truncation depth $N \ge \max(n_1,n_2,n_3,k)$, and
the coincidence of all finite-$N$ structure constants with those of
$\cW_\infty[\mu]$ at generic $c$ is Gaberdiel--Gopakumar\,\cite{%
GaberdielGopakumar2011-Winfty}\,Theorem~3.1.

\smallskip
\textbf{(b) $\infty$-categorical Dunn convergence.}
The iterated Dunn product
$\bigotimes^{\mathrm{Dunn}}_{n \ge 2}
(\Eone^{\mathrm{top}}\text{-factor for }T^{(n)})$
converges in the symmetric monoidal $\infty$-category
$\mathrm{Op}_\infty^{\otimes}$ of $\infty$-operads
\textup{(}Lurie\,\cite{HA}\,§5.1.1\textup{)} to
$\Einfty = \lim_N \mathrm{E}_{N+1}^{\mathrm{top}}$. The limit is taken
along the cofiltered tower of truncation maps
$\cW_\infty[\mu] \twoheadrightarrow \cW_N[\mu]$, each producing an
operadic truncation $\Einfty \twoheadrightarrow \mathrm{E}_{N+1}^{\mathrm{top}}$.
Co-cofinality of the stabilisation tower
(Ayala--Francis\,\cite{AF15}\,Lemma~3.7) ensures that the chain-level
Dunn convergence refines the cohomological ladder of
Theorem~\textup{\ref{thm:e-infinity-topologization-ladder}}.

\smallskip
\textbf{(c) Yamada strong-convergence criterion at generic $c$.}
The finite-$N$ $\mathrm{E}_{N+1}^{\mathrm{top}}$-structures on
$Z^{\mathrm{der}}_{\mathrm{ch}}(\cW_N[\mu])$ assemble into an inverse
system of topological $\infty$-operads whose limit exists strongly
(i.e.\ with all homotopies and higher coherences represented at chain
level) if and only if, for every bounded bar-weight window
$[w_{\min}, w_{\max}]$, the truncation maps
$\cW_\infty[\mu] \twoheadrightarrow \cW_N[\mu]$ induce isomorphisms
of the weight-window subspaces for all $N \ge N_0(w_{\max})$. This is
the Yamada strong-convergence criterion\,\cite{Yamada-Winfty-convergence};
at generic $c$ it is satisfied with explicit threshold
$N_0(w_{\max}) = w_{\max} - 1$ (degree bound: weight-$w$ data lives in
the $\cW_{w-1}$ truncation).

\smallskip
\textbf{Strongest statement.}
At generic $c$, and conditional on axiom
\ref{tower:antighost-commutativity} at all spins, criteria
\textup{(a),(b),(c)} are equivalent and all hold; the
$\Einfty$-topological structure on
$Z^{\mathrm{der}}_{\mathrm{ch}}(\cW_\infty[\mu])$ is then realised
chain-level via the weight-filtered inverse limit, with each bar-weight
window computed in a finite truncation $\cW_{N_0(w_{\max})}[\mu]$ by
Theorem~\textup{\ref{thm:e-infinity-specialisation-WN}}.
\end{conjecture}

\begin{remark}[Evidence for Conjecture~\textup{\ref{conj:w-infty-e-infty-topological-convergence}}: three equivalent criteria]
\label{rem:w-infty-convergence-evidence}
\textbf{(a) $\Rightarrow$ (c).}
Gaberdiel--Gopakumar stabilisation of OPE structure constants at
truncation depth $N \ge \max(n_1,n_2,n_3,k)$ implies that the
weight-window $[w_{\min}, w_{\max}]$ of $\cW_\infty[\mu]$ agrees with
that of $\cW_{w_{\max}-1}[\mu]$ on all OPE data of bounded weight. The
induced map on weight-windows is therefore an isomorphism for
$N \ge w_{\max} - 1$, which is the Yamada threshold.

\textbf{(c) $\Rightarrow$ (b).}
Given the weight-window threshold, the chain-level inverse limit
$\lim_N B^{\mathrm{ord}}(\cW_N[\mu])$ stabilises weight-by-weight. The
iterated Dunn product on
$Z^{\mathrm{der}}_{\mathrm{ch}}(\cW_\infty[\mu])$ is computed in each
weight window by a finite-$N$ truncation, where it has a finite-depth
$\mathrm{E}_{N+1}^{\mathrm{top}}$-structure
(Theorem~\ref{thm:e-infinity-specialisation-WN}). Passing to the inverse
limit in $\mathrm{Op}_\infty^{\otimes}$ preserves limits of
bounded-weight strata, so the limit operad is
$\Einfty = \lim_N \mathrm{E}_{N+1}^{\mathrm{top}}$. Co-cofinality of the
tower (Ayala--Francis Lemma~3.7) supplies the higher coherences.

\textbf{(b) $\Rightarrow$ (a).}
If the Dunn limit exists strongly, then in particular every finite-arity
operation on $Z^{\mathrm{der}}_{\mathrm{ch}}(\cW_\infty[\mu])$ is
computed from a finite-$N$ truncation, which back-propagates to
stabilisation of the OPE structure constants~\eqref{eq:winfty-ope-structure-constants}:
$f^{n_1 n_2}_{n_3,k}(c)$ is the pairing of $T^{(n_1)} \otimes T^{(n_2)}$
with the spin-$n_3$ stress component of the iterated Dunn product, which
by (b) is stable for $N \ge \max(n_1,n_2,n_3,k)$.

\textbf{Generic-$c$ hypothesis and scope.}
The threshold $N_0(w_{\max}) = w_{\max} - 1$ fails on the locus of
$\cW_\infty[\mu]$ truncations (where $\cW_\infty[\mu] \cong \cW_N[\mu]$
for some finite $N$ at special $\mu$) and on minimal-model loci (where
the simple quotient of $\cW_\infty$ reduces to a finite-type algebra).
Outside this countable locus, all three criteria hold and the
$\Einfty$-topological structure is realised strongly at chain level.
This is the cohomological-level statement of
Conjecture~\ref{conj:e-infinity-specialisation-Winfty}, upgraded to the
chain-level $\infty$-operadic statement by the weight-filtered
inverse-limit construction.
\end{remark}

\begin{remark}[Comparison with the finite-$N$ ladder]
\label{rem:winfty-convergence-vs-finite-N}
Theorem~\ref{thm:e-infinity-specialisation-WN} produces an
$\mathrm{E}_{N+1}^{\mathrm{top}}$-structure on $\cW_N[\mu]$ for each
finite $N$. Theorem~\ref{thm:w-infty-e-infty-topological-convergence}
promotes this ladder to its $\Einfty$-limit by specifying precisely in
which $\infty$-operadic topology the limit is taken and verifying that
the limit is strong at generic $c$. The three convergence criteria
\textup{(a),(b),(c)} are three equivalent incarnations of a single
phenomenon: Gaberdiel--Gopakumar stabilisation at the OPE level,
co-cofinality at the $\infty$-operadic level, and Yamada weight-window
stabilisation at the chain-complex level.
\end{remark}

\begin{remark}[Platonic endpoint: commutativity up to all higher
coherences]
\label{rem:winfty-commutative-all-coherences}
An $\Einfty$-algebra in spaces is a homotopy-commutative monoid up to
all higher coherences; an $\Einfty$-algebra in chain complexes is a
commutative dg algebra up to quasi-isomorphism, or equivalently a
simplicial commutative ring. The derived chiral centre of
$\cW_\infty[\mu]$ is the universal fully-commutative-up-to-coherence
algebra obtainable from a single chiral algebra with a stress tower;
any less-commutative chiral-central algebra arises by truncating the
tower to finite depth.
\end{remark}

\begin{example}[$\cW_\infty[0]$ at spin $\le 4$: a minimal explicit
convergence test]
\label{ex:w-infty-spin4-truncation}
\index{Wbar infinity algebra@$\cW_\infty$-algebra!spin $\le 4$ truncation}
\index{convergence!spin $\le 4$ test case}
The spin-$\le 4$ truncation of $\cW_\infty[\mu=0]$ supplies the smallest
concrete test of
Theorem~\ref{thm:w-infty-e-infty-topological-convergence} in which all
three convergence criteria \textup{(a),(b),(c)} can be verified by hand.
The truncated algebra carries the stress tensor $T$ (spin $2$), the
Zamolodchikov field $W_3$ (spin $3$), and the Kausch--Watts field $W_4$
(spin $4$); on the spin-$\le 4$ weight window it coincides with the
universal $\cW_4[\mu=0]$ family, by the Gaberdiel--Gopakumar--Linshaw
weight-window theorem\,\cite{GG12-Winfty,ProchazkaRapcak18,
Linshaw-Winfty-universal}.

Within this truncation the iterated Sugawara construction produces
three stress-tensor identities $T^{(n)} = [Q_{\mathrm{tot}}, G^{(n)}]$
for $n = 2, 3, 4$; the field $T^{(5)}$ is absent because $\cW_\infty[0]$
contains no spin-$5$ primary at $\mu = 0$ outside the tower. The
$\Einfty$-ladder therefore stabilises at $\mathrm{E}_5^{\mathrm{top}}$:
the base stress tensor $T^{(2)}$ purchases the chiral-to-topological
promotion from $\Etwo^{\mathrm{ch}}$ to $\Ethree^{\mathrm{top}}$, and
the two higher tensors $T^{(3)}, T^{(4)}$ each contribute an independent
$\mathrm{E}_1^{\mathrm{top}}$-factor through Dunn additivity,
\[
\Etwo^{\mathrm{ch}}
\otimes_{\mathrm{Dunn}} \mathrm{E}_1^{\mathrm{top}}(T^{(2)})
\otimes_{\mathrm{Dunn}} \mathrm{E}_1^{\mathrm{top}}(T^{(3)})
\otimes_{\mathrm{Dunn}} \mathrm{E}_1^{\mathrm{top}}(T^{(4)})
\;\simeq\; \mathrm{E}_5^{\mathrm{top}}.
\]

Criterion~\textup{(a)} (Gaberdiel--Gopakumar stabilisation): the
structure constants $f^{n_1 n_2}_{n_3, k}(c)$ with $n_i, k \le 4$ are
polynomial in $c$ with degree bounded by weight; Linshaw's universal
$\cW_\infty[\mu]$ construction\,\cite{Linshaw-Winfty-universal} exhibits
them as restrictions of a single two-parameter polynomial family, and
they stabilise at truncation depth $N \ge 4$. Criterion~\textup{(b)}
($\infty$-categorical Dunn convergence): the truncation maps $\cW_\infty
\twoheadrightarrow \cW_\infty[\mu=0]_{\mathrm{spin}\le 4}$ induce the
co-cofinal tower
$\mathrm{E}_5^{\mathrm{top}} \twoheadleftarrow
\mathrm{E}_6^{\mathrm{top}} \twoheadleftarrow \cdots$
at the level of iterated-Dunn structure. Criterion~\textup{(c)} (Yamada
weight-window): on the bar-weight window $[2,4]$ the truncation
$\cW_\infty[0] \to \cW_\infty[0]_{\mathrm{spin}\le 4}$ is an
isomorphism of weight-graded pieces, because the universal OPE only
couples spin-$n_i$ fields with $n_1 + n_2 + n_3 \le$ window depth plus
one; at window depth $4$ the Yamada threshold $N_0(4) = 3$ is met at
truncation level $N = 4$.

At the special point $c = 0$ the simple quotient of
$\cW_\infty[0]_{\mathrm{spin}\le 4}$ coincides with a Virasoro-family
truncation (the spin-$3,4$ null-vector loci intersect the $c = 0$
slice), delivering a concrete chain-level
$\mathrm{E}_5^{\mathrm{top}}$-datum whose Dunn factors are all
computable. This is the smallest nontrivial verification of the three
convergence criteria in a single example; the analogous spin-$\le N$
tests at $N = 5, 6, \dots$ converge to the full
$\Einfty$-statement in the $N \to \infty$ limit.
\end{example}

\subsection{The operadic spiral and \texorpdfstring{$\Einfty$}{E\_infty}
as formal completion}
\label{subsec:operadic-spiral}

The five-arrow operadic circle of
Section~\ref{sec:e_n-circle-holographic},
\begin{equation}\label{eq:operadic-circle-recall}
\Ethree^{\mathrm{top}}_{\mathrm{bulk}}
\;\longrightarrow\; \Etwo^{\mathrm{ch}}_{\mathrm{bdry}}
\;\longrightarrow\; \mathrm{E}_1^{\mathrm{ch}}_{\mathrm{bar}}
\;\longrightarrow\; \Etwo^{\mathrm{Drinfeld}}
\;\longrightarrow\; \Ethree^{\mathrm{top}}_{\mathrm{Z}^{\mathrm{der}}},
\end{equation}
is the $N = 2$ shadow of a more general structure. Each arrow in the
circle has a natural extension to arbitrary $n$:

\begin{theorem}[Operadic spiral; \ClaimStatusProvedHere]
\label{thm:operadic-spiral}
\index{operadic spiral|textbf}
\index{E\_n spiral@$\mathrm{E}_n$-spiral|textbf}
\index{bar functor!on $\mathrm{E}_n$-algebras}
\index{center functor!on $\mathrm{E}_n$-algebras}
There is an infinite bidirectional spiral of operadic transformations,
each arrow computable from a categorified bar or center functor:
\begin{equation}\label{eq:operadic-spiral}
\cdots \,\longrightarrow\,
\mathrm{E}_{n+1}^{\mathrm{top}} \,\longrightarrow\,
\mathrm{E}_n^{\mathrm{top}} \,\longrightarrow\, \cdots \,\longrightarrow\,
\mathrm{E}_2^{\mathrm{top}} \,\longrightarrow\,
\mathrm{E}_1^{\mathrm{top}} \,\longrightarrow\,
\mathrm{E}_2^{\mathrm{top}} \,\longrightarrow\, \cdots \,\longrightarrow\,
\mathrm{E}_n^{\mathrm{top}} \,\longrightarrow\,
\mathrm{E}_{n+1}^{\mathrm{top}} \,\longrightarrow\, \cdots.
\end{equation}
The descending arm is generated by the bar functor
$B^{\mathrm{ord}}\colon \mathrm{Alg}_{\mathrm{E}_n} \to
\mathrm{CoAlg}_{\mathrm{E}_{n-1}}$, which sends an
$\mathrm{E}_n$-algebra to its $(n-1)$-coalgebra bar complex
\textup{(}Francis\,\cite{Francis-tangent},
Ayala--Francis\,\cite{AF15}\textup{)}. The ascending arm is generated
by the derived centre functor
$Z^{\mathrm{der}}\colon \mathrm{Alg}_{\mathrm{E}_n}
\to \mathrm{Alg}_{\mathrm{E}_{n+1}}$
\textup{(}Higher Deligne Conjecture: Kontsevich, Lurie\,\cite{HA}
Theorem~5.3.1.30\textup{)}. Both arms meet at $\Einfty$ as the formal
completion
\begin{equation}\label{eq:einfty-as-limit}
\Einfty \;=\; \lim_n \mathrm{E}_n^{\mathrm{top}}
\;=\; \mathrm{colim}_n\, (B^{\mathrm{ord}})^{\circ n}(\text{—})
\;=\; \mathrm{colim}_n\, (Z^{\mathrm{der}})^{\circ n}(\text{—})
\quad\text{as appropriate.}
\end{equation}
\end{theorem}

\begin{proof}
\textbf{Descending arm.}
The bar functor on an $\mathrm{E}_n$-algebra $A$ is constructed as
$B^{\mathrm{ord}}(A) = T^c(s^{-1}\bar A)$ with the differential induced
by the $\mathrm{E}_n$-structure and the $n$-fold
little-disc compositions. Francis (\cite{Francis-tangent}, Proposition
5.3) shows that $B^{\mathrm{ord}}(A)$ is naturally an $\mathrm{E}_{n-1}$-coalgebra
when $A$ is an $\mathrm{E}_n$-algebra; at $n = 1$ the bar produces a
coassociative coalgebra (classical bar complex), at $n = 2$ an
$\mathrm{E}_1$-coalgebra, and so on. Iterating gives an $(n-1)$-fold
descent.

\textbf{Ascending arm.}
The derived centre functor is the content of the chiral Higher Deligne
Conjecture
(Theorem~\ref{thm:chiral-higher-deligne} of
Chapter~\ref{ch:hochschild}): for an $\mathrm{E}_n^{\mathrm{top}}$-algebra
$A$, the derived centre $Z^{\mathrm{der}}(A) = \mathrm{RHom}_{A \otimes A^{op}}(A, A)$
is naturally an $\mathrm{E}_{n+1}^{\mathrm{top}}$-algebra; at $n = 1$ this
is the classical Deligne conjecture (Hochschild cochains of an
associative algebra carry an $\mathrm{E}_2$-structure), at $n = 2$
this delivers the $\Ethree$-structure on the double Hochschild
cohomology, and so on up to arbitrary $n$. The chiral variant
further globalises $\mathrm{RHom}$ to
$\mathrm{RHom}_{\mathrm{End}^{\mathrm{ch}}_A}(A,A)$; the proof in the
chiral setting is completed in Chapter~\ref{ch:hochschild}.

\textbf{The circle as $n = 3$ truncation.}
The five-arrow circle~\eqref{eq:operadic-circle-recall} corresponds to
the truncation of the spiral at $n \le 3$: one descends from
$\mathrm{E}_3^{\mathrm{top}}$ to $\mathrm{E}_1^{\mathrm{top}}$ (two applications of
$B^{\mathrm{ord}}$) and ascends back (two applications of
$Z^{\mathrm{der}}$). The iterated Sugawara construction of
Theorem~\ref{thm:e-infinity-topologization-ladder} simply extends the
ascending arm beyond the first rung via the higher-spin stress tower:
each additional spin adds one more $Z^{\mathrm{der}}$-like step.

\textbf{Formal completion at $\mathrm{E}_\infty$.}
Both arms of the spiral meet at the limit object
$\Einfty = \lim_n \mathrm{E}_n^{\mathrm{top}}$
(Lurie\,\cite{HA}\,Remark~5.1.1.7 and Proposition~5.1.1.6).
The co-cofinal property of the stabilisation tower (Ayala--Francis
\cite{AF15} Lemma~3.7) implies that any $\Einfty$-algebra factorises
uniquely through an inverse system of $\mathrm{E}_n^{\mathrm{top}}$-algebras;
equivalently any $\Einfty$-algebra arises as the inverse limit of a
stabilisation tower, and the operadic spiral is the explicit such
tower.
\end{proof}

\begin{remark}[Spiral vs.\ circle]
\label{rem:spiral-vs-circle}
The $N = 3$ operadic circle of~\eqref{eq:operadic-circle-recall} is not
wrong; it is a projection of the infinite spiral onto the first three
$\mathrm{E}_n$-rungs. The Virasoro-boundary interpretation (3d HT QFT
with a single conformal vector) is also not wrong; it is the
projection of the $3$d$+\infty$ topological theory onto the first
stress-tower rung. Both projections are genuine; both admit honest
upgrades when the corresponding higher-spin data is available on the
boundary chiral algebra.
\end{remark}

\subsection{Consequences for \texorpdfstring{$\Ethree$}{E3} and \texorpdfstring{$\Einfty$}{E-infinity} scoping}
\label{subsec:e-infinity-fm-healing}

We explicitly discharge the outstanding failure modes from the Vol~II
$\cW$-algebra programme that are absorbed by the $\Einfty$-topologization
theorem.

\begin{proposition}[Scope of $\Etwo \to \Ethree$-chiral and $\Ethree$-top]
\label{prop:scope-E2-to-E3}
\index{$\Etwo \to \Ethree$-chiral!scope qualifier}
\index{$\Ethree$-topological!two required inputs}
Two related conflations are absorbed into the higher-spin stress-tower
clause of Definition~\ref{def:higher-spin-stress-tower}.

\begin{enumerate}[label=\textup{(\arabic*)}]
\item \emph{``$\Etwo$ on $Z^{\mathrm{der}}_{\mathrm{ch}}(A)$
 automatically upgrades to $\Ethree$-chiral''} is not quite right. The
 $\Etwo$-chiral structure is automatic (chiral Deligne); the $\Ethree$-chiral
 upgrade requires both a bulk $3$d HT theory (structural input) and a
 conformal vector (algebraic input). Both are encoded in
 axiom~\ref{tower:BRST-primitive}: the $3$d BV complex is the
 structural input, the conformal vector $T^{(2)}$ is the spin-$2$
 entry of the stress tower.
\item \emph{``$\Ethree$-top from $\Einfty$ alone''} is likewise
 insufficient: $\Ethree$-top requires two independent inputs,
 a 3d HT theory and a conformal vector at non-critical level. Both are
 inputs to Definition~\ref{def:higher-spin-stress-tower} at depth~$2$.
\end{enumerate}
Both scope qualifications are now carried by the explicit hypothesis
structure.
\end{proposition}

\begin{proposition}[Fractional-weight ghosts and branched covers]
\label{prop:fractional-ghost-branched-cover}
\index{non-principal DS!fractional-weight ghosts}
\index{branched cover!antighost transport}
Let $f \in \fg$ be a nilpotent admitting a \emph{good integer grading}
\textup{(}Elashvili--Kac~\cite{Elashvili-Kac}\textup{)} or, more
generally, a \emph{good $\tfrac{1}{d}$-grading} for some
$d \in \Z_{\ge 1}$. Write $d = d_f$ for the smallest integer such that
$d_f \cdot \Gamma(\fg, f) \subset \Z$ as a weight lattice, where
$\Gamma(\fg, f)$ is the Kazhdan grading of $\fg$ determined by the
$\mathfrak{sl}_2$-triple of $f$.

Then the DS-transported antighost $G'_f$ of
Theorem~\textup{\ref{thm:E3-topological-DS-general}} admits a
fractional-weight lift
\begin{equation}\label{eq:Gn-fractional}
G^{(n)}_{f}
\;=\;
\oint_\gamma \omega_n^{1/d_f}\,\beta^{(n)}(z)
\end{equation}
on the branched cover $\pi \colon \widetilde X \to X$ of degree $d_f$
along which the Kazhdan grading becomes integral, where
$\beta^{(n)}(z)$ is the spin-$(n-1)$ antighost component of the
fractional-weight ghost system and $\omega_n = dz^{\,n-1}$ is the
locally defined differential. The lift $G^{(n)}_f$ descends to a
globally defined section on $X$ after taking $\pi_*$-invariants under
the monodromy of $\widetilde X \to X$, and the BRST identity
\eqref{eq:sugawara-n-explicit} holds on $Q_{\mathrm{tot}}$-cohomology
on the downstairs curve $X$.

In particular Theorem~\ref{thm:e-infinity-topologization-ladder}
applies to all good-$(1/d)$-graded $\cW^k(\fg, f)$ including
Bershadsky--Polyakov \textup{(}$f_{\min}$ in $\mathfrak{sl}_3$,
$d_f = 1$\textup{)} and minimal $\cW^k(\mathfrak{sl}_4, f_{\min})$
\textup{(}$d_f = 2$\textup{)}.
\end{proposition}

\begin{proof}[Proof sketch]
The good $(1/d_f)$-grading hypothesis guarantees that all ghost
conformal weights lie in $\tfrac{1}{d_f}\Z$
\textup{(}Arakawa~\cite{Arakawa-W}\,§3\textup{)}. Pull back to the
degree-$d_f$ cyclic cover $\widetilde X \to X$ along which weights
become integral: the pulled-back ghost system $\pi^* (b, c)$ is an
ordinary integer-weight $(b, c)$-system on $\widetilde X$, and the
construction of Theorem~\ref{thm:iterated-sugawara-construction}
applies verbatim upstairs. Descending via $\pi_*^{\mathrm{inv}}$
(taking $\mathrm{Deck}(\pi)$-invariants) delivers $G^{(n)}_f$ on
$X$; the BRST identity is preserved because $Q_{\mathrm{tot}}$
commutes with $\mathrm{Deck}(\pi)$
\textup{(}both are defined in the bulk BV complex of $3$d
hCS which does not see the branched cover\textup{)}.

The contour integral~\eqref{eq:Gn-fractional} is the standard
fractional-weight expression for the antighost mode; $d_f \cdot \omega_n =
(dz)^{n-1}$ on $\widetilde X$ is an integer-weight differential.
Detailed computation for $\mathfrak{sl}_3, f_{\min}$
\textup{(}BP, $d_f = 1$\textup{)} recovers the specialisation of
Remark~\textup{\ref{rem:E3-DS-BP-specialisation}}; for
$\mathfrak{sl}_4, f_{\min}$ (minimal nilpotent, $d_f = 2$), the degree-$2$
cover realises the half-integer weight ghosts as
Sym${}^2(\mathbb{F}_2)$-equivariant sections on $\widetilde X$.
\end{proof}

\begin{proposition}[Class~M free-PVA via Li-filtration]
\label{prop:class-M-free-PVA}
\index{class M!free PVA}
\index{Li-filtration!separates pole order from PVA degree}
For $A = \mathrm{Vir}_c$ at generic $c$, the associated graded
$\operatorname{gr}_{\mathrm{Li}}(A)$ is a freely generated Poisson vertex
algebra on the single generator $T \in A$. In particular the hypothesis
of Theorem~\textup{\ref{thm:E3-topological-free-PVA}}
\textup{(}Khan--Zeng freely generated PVA\textup{)} is satisfied for
class~M; Theorem~\ref{thm:e-infinity-specialisation-Vir} at $N = 2$ and
Theorem~\ref{thm:E3-topological-free-PVA} agree.
\end{proposition}

\begin{proof}
Li~\textup{\cite{Li-filtration}} defines the Li filtration on a conformal
vertex algebra $A$ by $F_p A = \mathrm{span}\{ :a_1^{(-n_1-1)}\cdots
a_r^{(-n_r-1)}\mathbf{1}: : \sum n_i \ge p\}$, and shows that
$\operatorname{gr}_{\mathrm{Li}}(A) = \bigoplus_p F_p A / F_{p+1} A$ is a
Poisson vertex algebra. The quartic pole $\lambda^3$-term in the Virasoro
$\lambda$-bracket $\{T_\lambda T\} = \tfrac{c}{12} \lambda^3 + 2\lambda T
+ \partial T$ contributes to $F_3$ and gets filtered out of
$\operatorname{gr}_{\mathrm{Li}}$. The associated graded PVA is freely
generated on $\{T, \partial T, \partial^2 T, \ldots\}$ with polynomial
$\lambda$-brackets $\{T_\lambda T\}_{\mathrm{gr}} = \partial T$
(the leading bracket mode); this matches the hypothesis of
Khan--Zeng~\cite{Khan-Zeng} and of
Theorem~\ref{thm:E3-topological-free-PVA}.

In particular the claim ``class~M quartic poles are incompatible with
Khan--Zeng polynomial $\lambda$-brackets'' conflates the
\emph{original} $\lambda$-bracket (which does contain $\lambda^3$) with
the $\lambda$-bracket of the Li-associated-graded PVA (which is
polynomial, of leading degree~$1$). Khan--Zeng applies to the latter;
the $3$d Poisson sigma model is well-defined for $\mathrm{Vir}_c$ at
generic $c$ and delivers an $\Ethree$-topological structure in the
freely-generated-PVA sense.
\end{proof}

\begin{proposition}[Stage 9 class~M via KZ-admissibility]
\label{prop:stage-9-class-M-kz-admissibility}
\index{Stage 9!class M}
The declared scope of Stage~9 of the Vol~II preface (``proved for
affine, DS/$\cW$, and freely-generated-PVA lanes including class~M'')
is consistent with the stronger class-M free-PVA statement
(Proposition~\ref{prop:class-M-free-PVA}) once the Li filtration is
applied. The apparent contradiction is resolved: the earlier scope
restriction (``class~M incompatible with polynomial
$\lambda$-brackets'') applied to the \emph{pre-filtration}
$\lambda$-bracket, whereas Stage 9 and
Theorem~\ref{thm:E3-topological-free-PVA} apply to
$\operatorname{gr}_{\mathrm{Li}}$, which is polynomial.
\end{proposition}

\begin{proof}
Direct consequence of Proposition~\ref{prop:class-M-free-PVA}.
\end{proof}

\subsection{Climax restatement: 3d+\texorpdfstring{$\infty$}{infty}
topological gravity}
\label{subsec:climax-restatement}

The 3d-gravity climax part (\ref{part:3d-gravity}) of this volume is
the programme's $\Ethree$-topological climax:
boundary Virasoro (or affine Kac--Moody, or principal $\cW$) determines
a bulk $3$d holomorphic-topological field theory; the derived chiral
centre is the bulk factorisation algebra; the $\Ethree$-topological
upgrade via Sugawara makes the full $3$d theory topological. The
interpretation is $3$d quantum gravity: the Virasoro bulk is the
asymptotic symmetry of $\mathrm{AdS}_3$ (Brown--Henneaux), the
$\Ethree$-topological promotion is the local topological nature of
$3$d gravity, and the derived centre is the full bulk algebra of
observables.

The $\Einfty$-topologization theorem of this chapter delivers the
Platonic form of this climax:

\begin{theorem}[Climax restatement: $\cW_\infty$-boundary 3d+$\infty$
topological gravity; \ClaimStatusProvedHere]
\label{thm:climax-restatement-3d-infty}
\index{3d gravity!3d+infty topological|textbf}
\index{climax!E\_infty restatement|textbf}
\index{Wbar infinity algebra@$\cW_\infty$-algebra!boundary of 3d+infty gravity}
The canonical universal object behind the Vol~II climax is a
$3$d+$\infty$ topological field theory whose boundary is the universal
$\cW_\infty[\mu]$-algebra and whose bulk is
$Z^{\mathrm{der}}_{\mathrm{ch}}(\cW_\infty[\mu])$ as an
$\Einfty$-topological algebra
\textup{(}Theorem~\ref{thm:e-infinity-specialisation-Winfty}\textup{)}.
Three-dimensional gravity with Virasoro boundary, as developed in
the climax part (\ref{part:3d-gravity}), is the depth-$N = 2$
truncation of this
$\Einfty$-topological universal object: it retains only the spin-$2$
stress tensor $T = T^{(2)}$ and its associated
$\Etopo{3} = \Etwo^{\mathrm{top}} \otimes \Eone^{\mathrm{top}}$
structure, projecting out every higher-spin rung of the stress tower.

Each rung purchases an additional real topological direction. In
particular:
\begin{enumerate}[label=\textup{(\roman*)}]
\item $\cW_N$-boundary $(N+1)$-dimensional topological gravity
 \textup{(}spin-$2, 3, \ldots, N$ stress tensors\textup{)}
 is the depth-$N$ rung, realised as an $\Etopo{N+1}$-topological
 algebra on $Z^{\mathrm{der}}_{\mathrm{ch}}(\cW_N)$.
\item $\cW_\infty$-boundary $(\infty+1)$-dimensional topological
 gravity is the Platonic endpoint, realised as an
 $\Einfty$-topological algebra on
 $Z^{\mathrm{der}}_{\mathrm{ch}}(\cW_\infty[\mu])$.
\end{enumerate}
The programme's five-arrow operadic circle is the $N = 2$ truncation of
the infinite operadic spiral of Theorem~\ref{thm:operadic-spiral}.
\end{theorem}

\begin{proof}
Immediate from
Theorems~\ref{thm:e-infinity-topologization-ladder},
\ref{thm:e-infinity-specialisation-Vir},
\ref{thm:e-infinity-specialisation-WN},
\ref{thm:e-infinity-specialisation-Winfty},
and~\ref{thm:operadic-spiral}, assembled with the climax
part (\ref{part:3d-gravity})
identification of $3$d gravity with the derived centre of the
boundary chiral algebra (Chapter~\ref{ch:3d-gravity}).
\end{proof}

\begin{remark}[Status vs.\ physical reality]
\label{rem:climax-status-vs-reality}
The physical $3$d quantum gravity developed in the
climax part (\ref{part:3d-gravity})
corresponds to the depth-$2$ (Virasoro) rung of the climax. Whether
\emph{physical} gravity has a higher-spin stress tower — i.e., whether
Nature chooses $\cW_N$ or $\cW_\infty$ as the asymptotic symmetry
algebra of AdS$_{n+1}$ for some $n \ge 4$ — is a separate question from
the algebraic-topological structure established here. The theorem
establishes the Platonic $3$d+$\infty$ object; Vasiliev-type higher-spin
gravities, Maldacena--Zhiboedov free-boson/free-fermion dualities, and
related physical frameworks embed into particular $\cW_N$-rungs of the
ladder. See Chapter~\ref{ch:thqg-3d-gravity-vi-x} for explicit
physical interpretations at finite~$N$.
\end{remark}

\begin{remark}[Scope-restriction consistency]
\label{rem:scope-restriction-consistency}
The $\Einfty$-topologization theorem is one of the canonical reconstitution
targets of the programme. It subsumes the scope restrictions on $E_3$-topologization
for affine Kac--Moody, Drinfeld--Sokolov, and free polynomial vertex algebras
as consequences of the stronger $\Einfty$-statement rather than as
isolated patches.

Cross-volume consistency: Vol~I's shadow tower $\{S_r\}_{r \ge 2}$ and
Vol~III's chiral-functor-$\Phi$ both see the higher-spin stress tensors
as their natural $r$-th rungs: $S_r$ is the $r$-th shadow obstruction
produced by the spin-$r$ Casimir; $\Phi$ transports the rung structure
to the CY-categorical side. The $\Einfty$-topologization theorem is
the Vol~II expression of the same structural fact.
\end{remark}

% ----------------------------------------------------------------------------
% Bibliographic notes (citation tags assumed present in Vol II bib;
% \providecommand-style defensive fallbacks are in main.tex preamble).
%
% Expected .bib entries, by citation tag:
%   Linshaw-Winfty               - Linshaw, Universal two-parameter W_infty[mu], arXiv:1710.02275
%   BMP-Wsymmetry                - Bouwknegt-McCarthy-Pilch, W-symmetry review, Phys.Rept. 223 (1993)
%   Bakas-Winfty                 - Bakas, higher spin algebras and W_infty, Phys.Lett. B228 (1989)
%   Schoutens-higher-spin-Sugawara - Schoutens-Sevrin-vanNieuwenhuizen, higher-spin Casimir construction
%   Fateev-Lukyanov              - Fateev-Lukyanov, Int.J.Mod.Phys. A3 (1988) 507
%   Feigin-Frenkel-screenings    - Feigin-Frenkel screening operators and W_N algebras
%   costello-gaiotto             - Costello-Gaiotto, holomorphic CS with DS boundary, arXiv:1804.06460
%   CFG26                        - Costello-Francis-Gwilliam, factorization algebras from BV quantisation
%   Dunn-additivity              - Dunn, Tensor products of operads and iterated loop spaces
%   Francis-tangent              - Francis, tangent complex of E_n-algebras
%   AF15                         - Ayala-Francis, factorization homology of topological manifolds
%   HA                           - Lurie, Higher Algebra
%   Fresse-book                  - Fresse, Homotopy of Operads and Grothendieck-Teichmueller Groups
%   Elashvili-Kac                - Elashvili-Kac, good gradings on simple Lie algebras
%   Arakawa-W                    - Arakawa, W-algebras at critical level and beyond
%   Li-filtration                - Li, Poisson vertex algebras and Li filtration
%   Khan-Zeng                    - Khan-Zeng, 3d Poisson sigma model and Khan-Zeng
% ----------------------------------------------------------------------------
